\phase{3 --- ImageProcessingProtocol core implementation}

\textbf{Objective.} Implement the domain contract, typed composition, and declarative policy model against the generic execution dependency.

\deliverables
\begin{enumerate}
  \item Stable ImageIdentifier and document-scoped LayerIdentifier/LayerReference; monotonic ImageRevision; immutable ImageDocumentSnapshot with duplicate-layer and foreign-reference validation.
  \item Thread-safe reference-type Image state owner with synchronous immutable ImageStateSnapshot capture, a short lock-protected compare-and-swap commit, and atomic publication of document plus matching ImageStorageBindingSet.
  \item ImageExtent/ImageCanvasDescriptor and tagged LayerDescriptor content for raster, mask, label map, and annotation layers.
  \item LayerSpatialMapping: validated invertible 4x4 affine transform from layer sample-centre index coordinates to canvas sample-centre index coordinates; equal extents may omit it (identity), differing pixel extents must declare it.
  \item PixelDataDescriptor with component precision, channel identity/role/order, color or scientific interpretation, alpha mode, interleaved/planar layouts, plane subsampling, explicit sample-origin/chroma siting for subsampled planes, required base/row/slice alignment, checked stride/offset arithmetic, and minimum-storage-capacity validation.
  \item PixelPlaneLayout maps plane sample indices to layer-grid coordinates using sample origin plus subsampling factors; LayerSpatialMapping then maps layer-grid coordinates to canvas-grid coordinates.
  \item PixelStorageBinding maps every plane exactly once to one or more generic ExecutionResource handles, supports shared-allocation or per-plane allocation topology, checks capacity per assigned resource, and retains resource handles.
  \item ImageStorageBindingSet associates layer storage bindings with one exact ImageIdentifier/ImageRevision, verifies layer content/PixelDataDescriptor identity, and reports unbound pixel layers without rejecting non-resident document content.
  \item The generic associated-type implementation protocol and public specification descriptors.
  \item Typed sequential ImageFunction<Input,Output> composition, a zero-step identity function, and heterogeneous operation-step descriptors carrying immutable operation instances/specifications.
  \item Unique per-occurrence OperationStepIdentifier values; right-hand steps are reidentified during composition while source identifiers remain available for diagnostics.
  \item A fluent \texttt{qualifying} builder on \texttt{ImageFunction} that attaches an immutable \texttt{QualificationPolicy}; declarations retain scope and source provenance.
  \item Separate policy fields for hard requirements and optimization preferences; precision policies identify storage, computation, accumulation, or output precision. Storage dimensions accept PixelComponentType values; computation/accumulation dimensions accept ArithmeticPrecision values.
  \item A builder that validates provider/precision allowed-set intersections at effective pipeline/step scope, detects incompatible execution-capability value requirements, validates finite positive preference weights, and rejects contradictory same-objective preferences at one scope.
  \item Composition rules: pipeline-wide hard requirements apply to applicable steps; step policies remain scoped and are rebased when the right-hand function's step occurrences are reidentified; hard allowed sets combine by intersection; preferences combine only through explicit objective/priority/weight semantics.
  \item Applicability validation through OperationSpecification. Input-storage, computation, accumulation, and output precision applicability are declared independently. A step-scoped rule is rejected when that exact dimension is absent; pipeline rules/preferences for a dimension are rejected if no step exposes it. Deterministic-result applicability is declared separately.
  \item Explicit layer/composite targets and new-image/replacement result semantics.
\end{enumerate}

\exitcriteria
\begin{itemize}
  \item Identity/state tests prove layer identifiers are scoped by image identity, revisions never wrap, snapshots reject duplicate layer IDs, and foreign references cannot resolve.
  \item Image concurrency tests prove stable identity, immutable old snapshots, stale-revision rejection, validation failure atomicity, at-most-one concurrent commit for a shared expected revision, and resource retention by captured snapshots.
  \item Pixel descriptor tests cover UInt8/UInt16/UInt32, Int16/Int32, FP16/FP32/FP64, RGB/grayscale/YCbCr/custom color models, scientific channels/units, alpha, interleaved BGR, planar subsampled planes, centered versus co-sited chroma, base/row/slice alignment, volumetric strides, buffer-capacity checks, and arithmetic overflow.
  \item Layer-content tests prove masks, integer label maps, and annotations retain distinct semantics and validation.
  \item Spatial-registration tests reject mismatched layer/canvas extents without a mapping, accept explicit affine mappings, reject singular/non-affine/non-finite mappings, and verify coordinate transformation. Mapping does not perform resampling.
  \item Pixel storage-binding tests cover one shared allocation, independent per-plane allocations, capacity validation, complete/unique plane assignment, duplicate resource grouping, resource base-alignment compatibility, and retained handle lifetime.
  \item ImageStorageBindingSet tests reject foreign/unknown/annotation layers, duplicate bindings, descriptor mismatch, and stale revision association; partial sets report unbound pixel layers.
  \item Operation tests prove immutable parameters and versioned semantic identity.
  \item Composition enforces exact generic input/output type compatibility, preserves step order, and produces unique occurrence IDs when functions are reused.
  \item Identity composition leaves the sequence unchanged and no operation executes during composition.
  \item Builder tests reject incompatible required precision/provider declarations, empty allowed sets, incompatible execution-capability value kinds, invalid weights, unknown scopes, and contradictory objective declarations at construction/finalization.
  \item Composition tests reject empty intersections (including three-way intersections that are pairwise nonempty) and report every prior contributing source plus the incoming conflicting declaration and affected operation/step.
  \item Policy tests cover three-way allowed-set intersections (not only pairwise conflicts), execution-capability value-type conflicts, unknown step scopes, duplicate step IDs, invalid weights, and same-scope objective conflicts.
  \item Pipeline and step preferences for the same objective may use different priorities without leaking scope; differing priorities at the same scope are rejected.
  \item Right-hand step-scoped policy declarations are rebased to the new occurrence IDs when functions are composed.
  \item Precision-value validation rejects PixelComponentType values in computation/accumulation dimensions and ArithmeticPrecision values in inputStorage/output dimensions.
  \item Operation applicability validation treats inputStorage, computation, accumulation, and output precision independently; declaring one does not implicitly declare the others.
  \item A valid pipeline-level FP32 computation requirement is preserved through composition and is not silently weakened by a local preference.
  \item Public APIs expose no provider-native resource classes and no public node/graph editor model.
\end{itemize}

\textbf{Gate.} Do not add runtime routing until policy values are immutable, composable, explainable, and validated.

\phase{4 --- Environment discovery and implementation registries}

\textbf{Objective.} Connect registered implementations to the actual execution environment without making either protocol depend on concrete backends.

\deliverables
\begin{enumerate}
  \item Provider registry and discovery over the provider packages selected by the application. **Implemented:** an actor-backed registry rejects duplicate provider identities, provides deterministic registration snapshots, and discovers provider/device descriptors and state.
  \item Operation implementation registry that records operation identity/version, implementation identity/version, and declarative implementation specification. **Implemented:** an actor-backed registry validates registration, stores implementation versions independently, and enumerates candidates by operation or concrete step without ranking them.
  \item Internal type-erased registry entries that recover the concrete associated operation/input types before dispatch. **Implemented:** erased execution is module-internal and requires an explicit implementation key; public callers cannot bypass resolver qualification.
  \item Environment snapshot assembly from discovered provider/device descriptors, resource state, provider observations, and runtime-owned outstanding-work accounting. **Partially implemented:** provider/device snapshots and partial failures are preserved; runtime-owned outstanding-work accounting and snapshot lifecycle identity remain.
  \item Snapshot identity, observation time, freshness, and source-quality metadata.
  \item Runtime-owned outstanding-work ledger keyed by execution target. It distinguishes work reserved by the runtime from work accepted by a provider and releases submitted work only after completion is observed. **Implemented.**
  \item Identity-bearing runtime environment snapshots with a monotonic per-coordinator generation, captured time, provider discovery result, partial failures, and runtime-owned work-load snapshot. **Implemented.**
  \item Runtime workload counts remain separate from provider queue depth/utilization observations; missing provider telemetry remains unknown. **Implemented.**
  \item Registration/version compatibility checks and diagnostics for missing provider dependencies or unsupported operation versions.
  \item Test doubles for multiple implementations of one operation and for providers whose availability/load changes between snapshots.
\end{enumerate}

\exitcriteria
\begin{itemize}
  \item A consumer/sample target can register two implementation classes for the same operation without modifying \IPP{} or \IPR{} source. **Implemented/tested:** CPU and Metal test implementations share one GaussianBlur operation identity while declaring different static support.
  \item Registry registration validates implementation metadata against the associated operation contract; duplicate implementation identity/version pairs are rejected while separate versions coexist.
  \item Candidate enumeration is deterministic and does not rank or execute an implementation.
  \item The internal erased dispatch rejects mismatched operation/input concrete types even when stable type identifiers are reused.
  \item A registered implementation is not considered executable merely because its package loaded; its required provider and capabilities must be discovered.
  \item Environment snapshots distinguish capability from availability and distinguish unknown load from idle load.
  \item Runtime-owned reserved/submitted work counts remain separate from provider-reported queue depth and utilization.
  \item Each runtime snapshot has a unique identity and a monotonically increasing coordinator generation; the provider discovery and work-load snapshots expose their capture times.
  \item A submitted work lease remains counted until completion; cancellation requests do not release it prematurely.
  \item Heterogeneous registry dispatch is type-safe at the operation boundary and tested for mismatched operation/implementation pairs.
  \item Provider/implementation version incompatibilities produce explicit diagnostics.
  \item A provider whose state observation or device discovery fails does not erase other successful providers from the same environment snapshot; failures remain separately inspectable.
  \item A device whose state observation fails remains represented with unknown availability, not false availability or idle metrics.
  \item Duplicate device IDs and devices owned by another provider are rejected from the snapshot and reported as provider contract violations.
  \item A provider with no discovered devices can still represent an opaque managed target.
\end{itemize}

\textbf{Gate.} No routing decisions based only on platform name, package presence, or a static provider list.

\phase{5 --- Common qualification and resolution engine}

\textbf{Objective.} Implement the shared data-driven resolver that interprets operation semantics, implementation declarations, current environment state, and pipeline policy.

\deliverables
\begin{enumerate}
  \item One resolver entry point operating across all registered implementation classes; no mandatory custom qualification method on each implementation.
  \item Static eligibility evaluation over operation identity/version, validated parameters, input/output semantics, implementation support, provider/device capabilities, and hard pipeline constraints.
  \item Candidate-context preparation for statically eligible bindings, followed by context-scoped capability evaluation and verification of the corresponding optional context protocol conformance before any operation work is submitted.
  \item Structured result categories distinguishing semantically rejected, requirement-incompatible, temporarily unavailable, and eligible candidates.
  \item Ranking of eligible candidates according to pipeline preferences and the best reliable current environment observations.
  \item Initial cost estimates covering queue delay, transfers, conversions, allocation, compute, synchronization, and output materialization.
  \item Precision resolution across computation, accumulation, input/storage, and output precision. Required FP32 excludes candidates that cannot meet the specified FP32 dimension; preferred FP32 ranks eligible FP32 candidates higher without permitting semantic contract violations.
  \item Dynamic-state routing: a busy GPU may rank below a compatible CPU candidate for a latency-oriented request; an explicit required-provider constraint must prevent fallback.
  \item Candidate and decision diagnostics, including why each candidate was excluded, temporarily unavailable, or ranked lower.
  \item Stable resolver behavior for an identical operation, policy, registry, and environment snapshot; nondeterministic telemetry/model inputs must be recorded.
\end{enumerate}

\exitcriteria
\begin{itemize}
  \item The same GaussianBlur operation is evaluated against CPU and GPU implementation specifications by the same resolver.
  \item The resolver rejects unsupported border modes, precision, formats, or provider capabilities for the correct reason.
  \item With FP32 computation required, no FP16-only candidate is selected.
  \item With FP32 preferred, only candidates permitted by the operation's numerical contract can be ranked as alternatives.
  \item A context-scoped allocation/submission requirement is not satisfied by a provider/device declaration alone; the created context must report the capability and conform to the corresponding optional protocol.
  \item A busy/unavailable GPU routes to CPU only when CPU is registered, qualified, and fallback is permitted.
  \item A required CUDA request never silently falls back to CPU.
  \item An unavailable provider is reported as an unsatisfied environmental requirement, not as a contradictory pipeline policy.
\end{itemize}

\textbf{Gate.} No execution may begin from a candidate that has not passed the resolver's eligibility stage.

\phase{6 --- Execution coordination and pipeline planning}

\textbf{Objective.} Execute the resolver's choice safely and plan composed functions with data movement and resource lifecycle included.

\deliverables
\begin{enumerate}
  \item An execution coordinator that prepares, reserves/acquires, submits, awaits, validates, and releases execution resources through \GEP{}.
  \item Snapshot freshness and pre-execution revalidation. If a critical reservation fails, refresh and resolve again only within the original hard constraints.
  \item A first pipeline planner that considers operation adjacency, provider transitions, data residency, conversion requirements, and complete estimated cost.
  \item Explicit handling of managed runtimes whose internal scheduling is opaque; the outer runtime treats them as one implementation candidate unless reliable partition/resource facts are exposed.
  \item Cancellation, error propagation, device-loss handling, and resource cleanup.
  \item New-image result creation and transactional replacement: output is validated before publishing a new image revision.
  \item Diagnostics recording operation/implementation/provider/device, environment snapshot, candidate ranking, transfers, conversions, fallback, and execution outcome.
\end{enumerate}

\exitcriteria
\begin{itemize}
  \item An operation executes through the same path whether its implementation uses CPU, a GPU provider, or an opaque algorithm runtime.
  \item A sequence is not dispatched greedily without accounting for residency and transition costs.
  \item Resource reservation failure causes safe re-resolution or a structured failure; it never weakens hard requirements.
  \item Cancellation and device loss do not corrupt image state or release resources before work is complete.
  \item Failed execution or output validation leaves the original image revision unchanged.
  \item End-to-end tests prove that pipeline policy, resolver choice, and executed implementation agree.
\end{itemize}

\textbf{Gate.} Do not add optimization shortcuts until the baseline lifecycle is correct and observable.